% 译者 t06：ch04c（原書 159–168 ＝ 页图 p170–p179）
% 块边界说明：原書 p.158 末的“阶段 6”判定条件句已在 ch04b 内写完，本块自 p.159 的 [証明] 起。
% 定理 4.39 的证明跨原書 p.166–p.169：本块按块边界在 p.168 末闭合 proof（会多印一个 ■，原书此处无），
% ch04d 应以 [证明（续）] 从 p.169 首行接续。

% [原文 p.159]
\begin{proof}
假设 $a\ge 4$ 且 $b\ge 6$。令 $Y$ 由
\[
p^2+q^3\cdot u(qr^2,r)+qr^{a-4}\cdot u_a(qr^2,r)+r^{b-6}\cdot u_b(qr^2,r)=0
\]
定义。此时
\[
\pi\colon (p,q,r)\mapsto (x=p\cdot r^3,\quad y=q\cdot r^2,\quad z=r)
\]
确定双有理态射 $\pi\colon Y\to X$，而其不可约例外除子由 $(r=0)\subset Y$ 给出。在 $(r=0)$ 的一般点处 $Y$ 非奇异。$(1/x)(dy\wedge dz)$ 是 $\omega_X$ 的局部生成元，其拉回
\[
\pi^*\frac{dy\wedge dz}{x}=\frac1r\,dq\wedge dr
\]
沿 $r=0$ 有极点。因此 $X$ 不是标准奇点。
\end{proof}

\begin{proof}[定理 4.22 的证明（续）]
阶段 6 的逆命题通过考察 3 种情形得到，其中 2 种可以同时处理。

\noindent\textbf{阶段 7}\quad 设 $a\ge b-1$，$b=4,5$，并且
\[
f=y^3\cdot(\text{单位})+yz^a\cdot v_1+z^b\cdot(\text{单位}).
\]
若 $a\ge b$，则 $yz^a\cdot v_1$ 可被最后一项吸收，故由 4.24(4) 证明结束。于是设 $a=b-1$。对 $z^b$ 应用 4.24(2)，把 $y^iz^{b-i}$（$i\ge 3$）的倍数吸收进 $y^3\cdot(\text{单位})$，得到
\[
f=y^3\cdot(\text{单位})+y^2z^{b-2}\cdot v_2+z^b\cdot(\text{单位}).
\]
再对 $y^3$ 应用 4.24(2)；由于 $2(b-2)\ge b$，得到
\[
f=y^3\cdot(\text{单位})+z^b\cdot(\text{单位}),
\]
由 4.24(4) 得到 $E_6$ 或 $E_8$ 的方程。

\noindent\textbf{阶段 8}\quad 设 $b\ge 5$，并且
\[
f=y^3\cdot(\text{单位})+yz^3\cdot(\text{单位})+z^b\cdot v_3 .
\]
作变换 $y=y_1z_1$、$z=z_1$ 并爆破原点，得到
\[
\bar f=y_1^3\cdot(\text{单位})+y_1z_1\cdot(\text{单位})+z_1^{b-3}\cdot v_3 .
\]
此时 $\mult_0\bar f=2$，且 $\bar f_2$ 不是平方，故由 4.24(3) 知 $\bar f=0$ 可约。$f=0$ 的双有理变换 $\bar f=0$ 可约，故 $f=0$ 也同理。
% [原文 p.160]
由 $f_3=y^3\cdot(\text{常数})$，$f$ 的一个因子具有 $y\cdot(\text{单位})+z^2\cdot(\cdots)$ 的形状。把它取为新坐标 $y$，即可把 $f$ 化为如下形状：
\[
f=y\bigl(y^2\cdot(\text{单位})+yz^2\cdot v_4+z^3\cdot(\text{单位})\bigr).
\]
对 $z^3$ 应用 4.24(2) 得
\[
f=y\bigl(y^2\cdot(\text{单位})+z^3\cdot(\text{单位})\bigr),
\]
再应用 4.24(4) 得 $f=y^3+yz^3$，这是 $E_7$-型。

\noindent\textbf{阶段 9}\quad 设 $\mult_0f\ge 4$。
此时证明 $X$ 不是 Du Val 奇点。与阶段 6 同样，令
\[
Y:=\bigl(p^2+r^{-4}\cdot f(qr,r)\bigr),
\]
则
\[
\pi\colon (p,q,r)\mapsto (x=p\cdot r^2,\quad y=q\cdot r,\quad z=r)
\]
把 $Y$ 映入 $X$。与阶段 6 同样可知 $X$ 不是 Du Val 的。

最后，定理 4.22 的最后一部分由命题 4.10 以及已证明的部分推出。
\end{proof}

\section{Du Val 奇点的同时奇点消解}
\takeaway{本节研究 Du Val 奇点的平坦族的同时间奇点消解：先给出同时奇点消解的定义与存在性定理 4.28，再借助半普遍形变空间把问题归约到命题 4.31，并用有理曲面上的逐次爆破显式构造所需族。}

本节的目标是研究 Du Val 奇点的平坦族的同时间奇点消解。

\begin{definition}{4.26}
设 $f\colon X\to S$ 是概形或解析空间的态射。$f$ 的\textbf{同时奇点消解}（simultaneous resolution）是指交换图
\[
\begin{tikzcd}
\bar X \arrow[r,"p"] \arrow[d,"\bar f"'] & X \arrow[d,"f"] \\
S \arrow[r,equal] & S
\end{tikzcd}
\]
其中 $p$ 固有、$\bar f$ 光滑（即平坦且各纤维非奇异），并且对每个 $s\in S$，纤维之间的态射 $p_s\colon\bar X_s\to X_s$ 是双有理的。
\end{definition}

若 $S$ 是一个点，则 $f\colon X\to S$ 的同时奇点消解无非就是 $X$ 的奇点消解。在绝大多数情形，同时奇点消解不存在。

% [Tutor Note: 同时奇点消解要求同一族内所有纤维相容地一起消解，比逐纤维消解强得多；例 4.27 说明除 Du Val 外这几乎总是不可能，故定理 4.28 须先取有限覆盖。]

% [原文 p.161]
\begin{example}{4.27}
设 $f\colon X\to C$ 是到非奇异曲线的平坦态射，且存在点 $0\in C$ 使 $f$ 在 $C\setminus\{0\}$ 上光滑。下列情形中 $f$ 没有同时奇点消解。
\begin{enumerate}[label=(\arabic*),leftmargin=*]
\item $X_0$ 有奇点，而 $X$ 非奇异。
\item $X_0$ 是约化的奇异曲线。
\end{enumerate}
又，下列情形中不存在 $f$ 的、在 $C\setminus\{0\}$ 上成为同构的同时奇点消解。
\begin{enumerate}[label=(\arabic*),leftmargin=*,start=3]
\item $\dim X_0\ge 3$，且 $X_0$ 有孤立超曲面奇点。（由 [Gro68]，$X$ 是可分解的。）
\end{enumerate}
但若在 (3) 中不假定“在 $C\setminus\{0\}$ 上同构”，则有 S.~Kov\'a\v{c} 的如下反例。
\[
X=\Bigl(\sum_{i=1}^n x_iy_i=0\Bigr)\subset\mathbb{A}^{2n},
\qquad
f\colon X\to C=\mathbb{A}^1,\quad f(x,y)=y_n+\sum_{i=1}^n x_i .
\]
取此式时，$n\ge 3$ 蕴涵 $f\colon X\to C$ 满足 (3) 的条件。然而由 $X$ 的理想 $(x_1,\dots,x_n)$ 作成的爆破 $p\colon\bar X\to X$ 是 $f$ 的同时奇点消解。
\end{example}

\begin{theorem*}[{[Bri71, Tyu70]}]{4.28}
设 $f\colon (x\in X)\to(0\in S)$ 是解析流形的芽之间的态射，$X_0$ 是在 $x$ 处有 Du Val 奇点的曲面。此时可以选取有限满射 $g\colon S'\to S$，使拉回 $f'\colon X':=X\times_S S'\to S'$ 具有同时奇点消解
\[
\begin{tikzcd}
\bar X' \arrow[r,"p"] \arrow[d,"\bar f'"'] & X' \arrow[d,"f'"] \\
S' \arrow[r,equal] & S' .
\end{tikzcd}
\]
进一步可以要求 $p$ 是射影的，且对每个 $s'\in S'$，$\bar X'_{s'}$ 是 $X'_{s'}$ 的极小奇点消解。这样的 $p$ 称为\textbf{射影的同时极小奇点消解}。
\end{theorem*}

% [原文 p.162]
本定理已知有几个证明。[Bri66] 的最初证明利用 Du Val 奇点全部形变的明确描述来构造同时奇点消解。Du Val 奇点与半单 Lie 群相联系，由幂零元集合的研究可得到定理 4.28 的另一个证明 [Bri71]。[Art74] 的证明使用了“同时奇点消解的栈（stack）”这一般概念。这一手法最为一般，但含有技术上的难解之处。这里给出仿照 [Tyu70]、使用具体的同时间奇点消解的证明：从奇点消解出发，通过收缩反过来造出奇点。这一方法在特殊情形也见于 [Pin80]。
此处的构造使用 Du Val 奇点的半普遍形变空间，这一概念将在 \S4.5 中详细讨论。

\begin{definition}{4.29}
设 $0\in X_0=(f(x_1,\dots,x_m)=0)\subset\C^m$ 是原点处孤立超曲面奇点的芽。由局部零点定理，
\[
\C[[x_1,\dots,x_m]]/(f,\partial f/\partial x_1,\dots,\partial f/\partial x_m)
\]
是有限维 $\C$-向量空间。其维数称为 $X_0$ 的 \textbf{Tyurina 数}，记作 $\tau(X_0)$。
\end{definition}

% [Tutor Note: 此处并非笔误：日文作“局所零点定理”，英文底本同处作 local Nullstellensatz；Nullstellensatz 中文常译“零点定理”（亦译“零化定理”），故照原书字面保留“局部零点定理”。]

\begin{example}{4.30}
通过具体计算可知 $\tau(A_n)=n$、$\tau(D_n)=n$、$\tau(E_n)=n$ 成立。
\end{example}

下列结果表明，为证明定理 4.28，只需对一个充分大的形变族构造同时奇点消解。

\begin{propositionp}{4.31}
设 $0\in X_0\subset\C^3$ 是 Du Val 奇点。假设存在解析流形的芽之间的平坦态射 $\bar f\colon (x\in\bar X)\to(0\in\bar S)$，具有以下性质。
\begin{enumerate}[label=(\arabic*),leftmargin=*]
\item $\bar X_0\cong X_0$。
\item $\dim\bar S=\tau(X_0)$。
\item $\bar X_0$ 是唯一具有同构于 $X_0$ 的奇点的纤维。
\item $\bar f$ 具有射影的同时极小奇点消解。
\end{enumerate}
此时定理 4.28 对 $X_0$ 成立。
\end{propositionp}

\begin{proof}
设 $u\colon (x\in\mathbf{X})\to(0\in\mathbf{U})$ 是 $X_0$ 的半普遍形变空间（定义 4.59）。
% [原文 p.163]
由定义存在态射 $u(\bar f)\colon\bar S\to\mathbf{U}$。命题 4.31(3) 意味着 $u(\bar f)^{-1}(0)=\{0\}$，故 $u(\bar f)$ 是有限态射。（例如参看 [GR84, p.63]。）$\bar S$ 与 $\mathbf{U}$ 同维数，且如定理 4.61 的具体描述所示 $\mathbf{U}$ 不可约。因此 $u(\bar f)$ 是满射。
设 $f\colon X\to S$ 是 $X_0$ 的任意平坦形变，则得到态射 $u(f)\colon S\to\mathbf{U}$。令 $S':=\bar S\times_{\mathbf{U}}S$，并设 $f'\colon X'\to S'$ 是由此诱导的形变族。$S'\to S$ 是满射，而 $f'\colon X'\to S'$ 的射影同时极小奇点消解可由 $\bar f\colon (x\in\bar X)\to(0\in\bar S)$ 的相应消解拉回得到。
\[
\begin{tikzcd}[row sep=large,column sep=large]
 & X' \arrow[rr] \arrow[dl] \arrow[dd] & & \bar X \arrow[dl] \arrow[dd,"\bar f"] \\
X \arrow[rr] \arrow[dd,"f"'] & & \mathbf{X} \arrow[dd,"u"] & \\
 & S' \arrow[dl] \arrow[rr] & & \bar S \arrow[dl] \\
S \arrow[rr,"u(f)"'] & & \mathbf{U} &
\end{tikzcd}
\]
\end{proof}

至此，为证明定理 4.28，对每个 Du Val 奇点只要构成命题 4.31 那样的例子便足够。此处给出的构成是全局的例子，所有奇点都在有理曲面上实现。

% [Tutor Note: 上图的方块把“半普遍族”的拉回与“同时消解”的拉回叠在一起：竖直方向的四条箭头是各族的底空间与总空间投影，斜箭头表示基变换。]

\par\medskip\noindent\textbf{构成 4.32}\quad 设 $Z$ 是非奇异曲面，$C\subset Z$ 是非奇异曲线。设 $P_1,\dots,P_n\in C$ 是（允许重复的）闭点集。如下定义曲面列。
\begin{enumerate}[label=(\arabic*),leftmargin=*]
\item $Z_0=Z$，$r_0=\mathrm{id}$。
\item 设已定义双有理态射 $r_i\colon Z_i\to Z_0$。由 $(r_i)_*^{-1}C\cong C$，可把点 $P_{i+1}$ 与 $Z_i$ 上的点 $\bar P_{i+1}$ 等同。令 $Z_{i+1}\to Z_i$ 为沿 $\bar P_{i+1}$ 的爆破，$r_{i+1}\colon Z_{i+1}\to Z_i\to Z_0$ 为合成态射。\hfill$\square$
\end{enumerate}

\par\medskip\noindent\textbf{注意}\quad 若 $Z$ 是曲面、$P_1,\dots,P_n\in Z$ 是闭点集，则一般而言不能定义 $Z$ 在这些点处的爆破。即使只有 2 个点，第 2 次爆破的选择也不唯一。若采用此处给出的构成，则可以处理由 $Z$ 逐次爆破得到的所有曲面的空间。像上面那样指定曲线 $C$，几乎可以消除技术上的麻烦问题。

% [原文 p.164]
这一构成具有如下性质是显然的。

\begin{lemma*}{4.33}
记号同上。此时下列结论成立。
\begin{enumerate}[label=(\arabic*),leftmargin=*]
\item $r_n\colon Z_n\to Z_0$ 与点 $P_i$ 的顺序无关。
\item $r_n$ 的例外曲线是 $(-1)$-曲线或 $(-2)$-曲线。
\end{enumerate}
\end{lemma*}

若使点 $P_i$ 变动，则所得曲面构成族。

\par\medskip\noindent\textbf{构成 4.34}\quad 设 $Z$ 是非奇异曲面，$C\subset Z$ 是非奇异曲线。对每个 $n\ge 0$，如下定义流形与态射的列。
\begin{enumerate}[label=(\arabic*),leftmargin=*]
\item $S=C^n$，$Y_0=Z\times S$，$p_0\colon Y_0\to S$ 为射影。$D_0\subset Y_0$ 是除子 $D_0=C\times S\subset Z\times S$。
\item $s_i\colon C^n\to Y_0$ 是截面 $s_i(P_1,\dots,P_n)=P_i\times(P_1,\dots,P_n)$，令 $B_i=\im(s_i)$。（$B_i$ 是 $C^n\times C\subset C^n\times Z$ 的“第 $i$ 条对角线”。）$B_i\subset D_0$，$\sum B_i\subset D_0$ 是 $D_0$ 的 snc 除子。
\item 设已定义双有理态射 $r_i\colon Y_i\to Y_0$。令 $D_i:=(r_i)_*^{-1}D_0$，并设 $r_i\colon D_i\to D_0$ 是同构。此时 $B_j$ 与部分概形 $\bar B_j\subset D_i$ 等同。令 $Y_{i+1}\to Y_i$ 为沿 $B_{i+1}$ 的爆破，$r_{i+1}\colon Y_{i+1}\to Y_i\to Y_0$ 为合成态射，$p_{i+1}=p_0\circ r_{i+1}$。此时注意 $(r_{i+1})_*^{-1}D_0\cong D_0$。\hfill$\square$
\end{enumerate}

$Y_n$ 依赖于爆破的顺序，但 $p_n$ 的纤维不依赖。

\begin{lemma}{4.35}
记号同上。此时 $(P_1,\dots,P_n)$ 上 $Y_n\to C^n$ 的纤维是构成 4.32 中定义的 $Z$ 在 $P_1,\dots,P_n\in C$ 处的爆破。
\end{lemma}

通过适当选取 $C\subset Z$，可得到各种 Du Val 奇点的同时奇点消解。

\begin{example}{4.36}
（$A_n$-型）设 $0\in\PP^2$ 是一点，$Z=B_0\PP^2$，$q\colon B_0\PP^2\to\PP^2$ 为射影。设 $C\subset Z$ 是过 $0$ 的直线的双有理变换，$P_0\in C$ 是 $C$ 与 $q$ 的例外曲线的交点。像构成 4.34 那样构成 $p_n\colon Y_n\to S$。令 $V:=\PP^2$、$v:=q\circ r_n\colon Y_n\to Y_0\to V$，以及 $L_r:=v^*\cO_V(r)(D_n)$。
\end{example}

\begin{example}{4.37}
（$D_n$-型）设 $Z=\PP^1\times\PP^1$，$q\colon Z\to\PP^1$ 为第 2 射影。设 $C\subset Z$ 是满足 $C\sim 2\cdot(\text{1 点})\times\PP^1+\PP^1\times(\text{1 点})$ 的非奇异曲线，$P_0\in C$ 是 $q$ 的分歧点。
% [原文 p.165]
像构成 4.34 那样构成 $p_n\colon Y_n\to S$。令 $V:=\PP^1$、$v:=q\circ r_n\colon Y_n\to Y_0\to V$，以及 $L_r:=v^*\cO_V(r)(D_n)$。
\end{example}

\begin{example}{4.38}
（$E_n$-型）设 $Z=\PP^2$，$C\subset Z$ 是非奇异 3 次曲线，$P_0\in C$ 是其拐点（flex）。像构成 4.34 那样构成 $p_n\colon Y_n\to S$。取 $q\colon Y_0\to V$（$V$ 为一点），令 $v:=q\circ r_n\colon Y_n\to Y_0\to V$，以及 $L_r:=v^*\cO_V(r)(D_n)$。（$L_r$ 与 $r$ 无关。）
\end{example}

\begin{theorem*}{4.39}
上述构成具有下列性质。
\begin{enumerate}[label=(\arabic*),leftmargin=*]
\item $p_n\colon Y_n\to S$ 非奇异且射影。
\item 对 $i=1,2$、$s\ge 1$、$r\gg 1$ 有 $R^i(p_n)_*L_r^s=0$。
\item 对 $s\gg 1$ 与 $r\gg 1$，$L_r^s$ 是 $p_n$-自由。
\item 由此确定态射 $F_r\colon Y_n\to\bar Y_n$，并得到如下交换图。
\[
\begin{tikzcd}[column sep=large]
Y_n \arrow[r,"F_r"] \arrow[d,"p_n"'] & \bar Y_n \arrow[r,"v"] \arrow[d,"\bar p_n"'] & V\times C^n \arrow[d] \\
S \arrow[r,equal] & S \arrow[r,equal] & C^n .
\end{tikzcd}
\]
\item $\bar p_n\colon\bar Y_n\to S$ 平坦、射影，且其任意纤维是只含 Du Val 奇点的曲面。
\item $p_n\colon Y_n\to S$ 是 $\bar p_n\colon\bar Y_n\to C^n$ 的射影同时极小奇点消解。
\item $A_n$-型：$\bar p_n$ 的纤维中具有 $A_n$-奇点的只是 $(P_0,\dots,P_0)\in S$ 上的纤维。
\item $D_n$-型（$n\ge 4$）：$\bar p_n$ 的纤维中具有 $D_n$-奇点的只是 $(P_0,\dots,P_0)\in S$ 上的纤维。其中 $P_0$ 是 $q\colon C\to\PP^1$ 的任意分歧点。
\item $E_n$-型（$n=6,7,8$）：若适当选取 $(P_0,\dots,P_0)\in S^0\subset S$ 的邻域，则在 $S^0$ 上 $\bar p_n$ 的纤维中具有 $E_n$-奇点的只是 $(P_0,\dots,P_0)$ 上的纤维。
\item 令 $Z_n(0):=p_n^{-1}(P_0,\dots,P_0)$、$\bar Z_n(0):=\bar p_n^{-1}(P_0,\dots,P_0)$。设 $\Gamma\subset Z_n(0)$ 是 $F_r\colon Z_n(0)\to\bar Z_n(0)$ 的例外曲线。此时 $\Pic Y_n\to\Pic\Gamma$ 是满射。
\end{enumerate}
\end{theorem*}

% [Tutor Note: 定理 4.39 的 (7)–(9) 是关键的刚性条件：无论爆破中心如何变动，纤维族中只有特定的中心配置才退化出指定的 Du Val 奇点。]

% [原文 p.166]
\begin{remark}{4.40}
(1) $Z_n(0)$ 上曲线的形状由下列图表示。
\begin{center}
\begin{tikzpicture}[scale=0.9]
\node at (-1.0,0.02) {$A_n$:};
\draw[thin] (-0.4,0)--(0.7,0);
\draw[thin] (0.7,0)--(1.9,0);
\draw[thin] (2.5,0)--(3.6,0);
\draw[thin] (3.6,0)--(4.6,0);
\draw[thin] (4.6,0)--(5.8,0);
\draw[thin] (0.7,0.07)--(0.77,0)--(0.7,-0.07)--(0.63,0)--cycle;
\draw[thin,fill=white] (1.9,0) circle (2pt);
\draw[thin,fill=white] (3.6,0) circle (2pt);
\fill (4.6,0) circle (2pt);
\node at (5.95,0) {$C'$};
\node at (2.2,0) {$\cdots\cdots$};
\end{tikzpicture}
\end{center}
\begin{center}
\begin{tikzpicture}[scale=0.9]
\node at (-1.0,0.02) {$D_n$:};
\draw[thin] (-0.4,0)--(0.7,0);
\draw[thin] (0.7,0)--(1.9,0);
\draw[thin] (2.5,0)--(3.6,0);
\draw[thin] (3.6,0)--(4.6,0);
\draw[thin] (4.6,0)--(5.6,0);
\draw[thin] (5.6,0)--(6.8,0);
\draw[thin,fill=white] (-0.4,0) circle (2pt);
\draw[thin,fill=white] (0.7,0) circle (2pt);
\draw[thin,fill=white] (3.6,0) circle (2pt);
\draw[thin,fill=white] (4.6,0) circle (2pt);
\draw[thin,fill=white] (5.6,0) circle (2pt);
\fill (6.8,0) circle (2pt);
\draw[thin] (0.7,0)--(0.7,-1.0);
\draw[thin] (0.7,-0.93)--(0.77,-1.0)--(0.7,-1.07)--(0.63,-1.0)--cycle;
\node at (7.15,0) {$C'$};
\node at (2.2,0) {$\cdots\cdots$};
\end{tikzpicture}
\end{center}
\begin{center}
\begin{tikzpicture}[scale=0.9]
\node at (-1.0,0.02) {$E_n$:};
\draw[thin] (-0.4,0)--(0.7,0);
\draw[thin] (0.7,0)--(1.9,0);
\draw[thin] (1.9,0)--(3.0,0);
\draw[thin] (3.6,0)--(4.8,0);
\draw[thin] (4.8,0)--(5.8,0);
\draw[thin] (5.8,0)--(6.9,0);
\draw[thin,fill=white] (-0.4,0) circle (2pt);
\draw[thin,fill=white] (0.7,0) circle (2pt);
\draw[thin,fill=white] (1.9,0) circle (2pt);
\draw[thin,fill=white] (4.8,0) circle (2pt);
\draw[thin,fill=white] (5.8,0) circle (2pt);
\fill (6.9,0) circle (2pt);
\draw[thin] (1.9,0)--(1.9,-1.0);
\draw[thin] (1.9,-0.93)--(1.97,-1.0)--(1.9,-1.07)--(1.83,-1.0)--cycle;
\node at (7.25,0) {$C'$};
\node at (3.1,0) {$\cdots\cdots$};
\end{tikzpicture}
\end{center}
图中 $\bullet$ 是最后一次爆破得到的 $(-1)$-曲线，$C'=(r_n)_*^{-1}C$，而 $\diamond$ 是 $(-2)$-曲线 $(r_n)_*^{-1}G$。其中 $G\subset Z$ 在 $A_n$-型时为例外曲线，$D_n$-型时为 $\PP^1\times q(P_0)$，$E_n$-型时为在 $P_0$ 处与 $C$ 相切的直线。
(2) 即使 $E_n$-型，$S$ 上 $\bar p_n$ 的纤维中具有 $E_n$-奇点的也只限于 $(P,\dots,P)$ 上的纤维（其中 $P$ 是拐点）。仔细分析状况即可得到这一点。考察纤维中可能出现的 Du Val 奇点的形状是非常有趣的。对 $E_6$，其答案对应于关于奇异 3 次曲线的古典结果（参看 [Hen11, Fur86]）。
\end{remark}

\begin{proof}
(1) 由引理 4.35 得到。下面内容是“上同调与基变换”定理（cohomology and base change）[Har77, III.12.11] 的特殊情形。

\par\smallskip\noindent\hspace*{2em}（上同调与基变换定理）$f\colon X\to Y$ 是 Noether 概形的平坦射影态射，$L$ 是 $X$ 上的可逆层，$y\in Y$ 是满足 $H^i(X_y,L_y)=0$（$\forall i>0$）的点。此时在 $y$ 的某个 Zariski 邻域 $U$ 上下列结论成立。
\begin{enumerate}[label=(\arabic*),leftmargin=2em]
\item $R^if_*L|_U=0$（$\forall i>0$），且 $f_*L|_U$ 局部自由；
\item $f_*L\to H^0(X_u,L_u)$ 对一切 $u\in U$ 是满射。
\end{enumerate}

由此可知，为证明 (2),(3),(4)，只需对每个纤维证明即可。设 $C\subset Z$ 是例中的一个。对任意 $P_1,\dots,P_n\in C$ 得到 $r_n\colon Z_n\to Z$ 与 $v\colon Z_n\to Z\to V$。$D_n$ 在 $Z_n$ 上的限制是曲线 $C_n=(r_n)_*^{-1}(C)$。对 $r\ge 0$，定义 $\Q$-除子 $M_r=r\cdot v^*\cO_V(1)+C_n$。对充分大的整数 $r$，可确认 $M_r$ 满足消没定理与无基点定理
% [原文 p.167]
的条件。
首先，$M_r$ 有效，且若 $(M_r\cdot C_n)>0$ 则 $M_r$ 是 nef 且大的，这是显然的。注意
\[
\begin{aligned}
(M_r\cdot C_n)&=r(r_n^*q^*\cO_V(1)\cdot C_n)+(C_n^2)\\
&=r(q^*\cO(1)\cdot C)+(C^2)-n .
\end{aligned}
\]
由此，若 $r\gg 1$ 则 $(M_r\cdot C_n)>0$。这是因为 $A,D$-型时 $(q^*\cO(1)\cdot C)>0$，$E$-型时 $(q^*\cO(1)\cdot C)=0$、$(C^2)-n=9-n>0$，都容易看出。因此由上所指出的，$r\gg 1$ 时 $M_r$ 是 nef 且大的。
由构成，在 $A,D,E$ 的任一情形都存在适当的 $t$，使 $Z$ 上 $C\sim-K_Z+t\cdot q^*\cO(1)$ 成立。由此 $C\sim-K_{Z_n}+tv^*\cO(1)$，进而
\[
sM_r-K_{Z_n}\sim(s+1)\Bigl(\frac{sr-t}{s+1}v^*\cO(1)+C_n\Bigr)=(s+1)M_{\frac{sr-t}{s+1}}
\]
成立。因此，若 $s>0$ 且 $r\gg 1$，则 $sM_r-K_{Z_n}$ 是 nef 且大的。这正是所应证明的消没定理、无基点定理的条件。
以下虽是记号的滥用，但 $r\in\Z_{>0}$ 时也把 $M_r$ 用作线丛。由消没定理 2.70，对 $i,s>0$、$r\gg 1$ 有 $H^i(Z_n,M_r^s)=0$。由无基点定理 3.3，对 $r,s\gg 1$，$M_r^s$ 由整体截面生成。
$\bar p_n$ 的纤维是由线性系 $|M_r^s|$ 给出的曲面 $Z_n$ 的像。取 $s\gg 1$，可以假定
\[
H^0(\bar Z_n,(\bar M_r^s)^t)=H^0(Z_n,(M_r^{st})),\quad \forall t\ge 1 .
\]
于是 $\bar p_n$ 的纤维的 Hilbert 多项式与 Hilbert 多项式 $\chi(Z_n,M_r^{st})$ 一致。因此由下列结果，$\bar p_n$ 平坦。

\par\smallskip\noindent\hspace*{2em}（平坦性判定法 [Har77, III.9.9]）$f\colon X\to Y$ 是 Noether 概形的射影态射，$Y$ 约化且不可约，$L$ 是 $f$-丰富的 Cartier 除子。此时 $f$ 平坦 $\iff$ 关于 $L_y$ 的 $X_y$ 的 Hilbert 多项式与 $y\in Y$ 无关。

$K_{Y_n}$ 在 $F_r$ 的纤维上数值平凡，故 $\bar p_n$ 的纤维只含 Du Val 奇点。由此得到 (5)，进而得到 (6)。

% [原文 p.168]
为使 $Z_n\to\PP^2=V$ 的例外集中含有 $A_n$-图，例外曲线必须含有由 $(n+1)$ 个不可约成分构成的连通曲线。（至少需要一条相交的 $(-1)$-曲线。）为此，所有爆破的中心必须位于 $0\in\PP^2$ 的原像内。这样得到 (7)。
$D$-型的情形同样，所有爆破的中心必须位于 $v$ 的同一纤维 $F$ 上的点。若 $C\cap F$ 是 2 点，则成为 $A_{n-1}$-图，故得到 (8)。
$E$-型的分析更难。设 $P_1(t),\dots,P_n(t)\in C$ 是依赖于参数 $t\in\C$ 的 $C$ 的点构成的解析族，且满足 $P_i(0)=P_0$。
\[
\PP^2=Z\xleftarrow{r_n(t)}Z_n(t)\to\bar Z_n(t)
\]
是对应的曲面族。假定对一切 $t$，$\bar Z_n(t)$ 有 $E_n$-奇点 $Q(t)\in\bar Z_n(t)$。此时必须证明对一切 $i,t$ 有 $P_i(t)=P_0$。
$Z_n(t)$ 的 Picard 数是 $n+1$，$\bar Z_n$ 的 Picard 数至少是 1，故 $Z_n(t)\to\bar Z_n(t)$ 至多只有 $n$ 条例外曲线（引理 3.40）。特别地，$Q(t)$ 是 $\bar Z_n(t)$ 的唯一奇点，且 $t\mapsto Q(t)$ 是解析的。这意味着族 $\bar Z_n(t):t\in\C$ 可由反复沿截面爆破得到。因此，适当编号 $Z_n(t)\to\bar Z_n(t)$ 的例外曲线 $J_1(t),\dots,J_n(t)$ 后，可使各 $J_i(t)$ 解析地依赖于 $t$。
具体计算可知，即便所有点 $P_1,\dots,P_n$ 一致，得到 $E_n$-图的也只限于 $P_1$ 为拐点之时。拐点的切线的双有理变换，以及 $r_n(0)$ 的例外曲线中的 $(n-1)$ 条，构成 $E_n$-图。可以假定 $J_1(0)$ 是拐点的切线的双有理变换，而 $J_2(0),\dots,J_n(0)$ 是例外的。于是 $r_n(t)(J_1(t))$ 是 $Z$ 内的直线，$J_2(t),\dots,J_n(t)$ 是 $r_n(t)$-例外的。因此 $r_n(t)\colon Z_n(t)\to Z(t)$ 只有唯一的例外 $(-1)$-曲线，且所有点 $P_1(t),\dots,P_n(t)$ 一致。由上述计算，各 $P_i(t)$ 也是拐点，故恒等地有 $P_i(t)=P_0$。
为证明 (10)，注意在 $A_n$（或 $D_n,E_n$）的情形，$\Gamma\subset Z_n(0)$ 是 $A_n$（或 $D_n,E_n$）-型的图。$Z_n(0)$ 的 Picard 群由 $Z$ 的 Picard 群与例外曲线生成。
% [Tutor Note: 证明 (10) 时把 Picard 群的计算归约到例外图 Γ 上：要在例外纤维上比较可逆层，只需在例外曲线上做除子计算即可。]
\end{proof}

% [Tutor Note: 原书此处印刷作“$T\mapsto Q(t)$ 是解析的”，英文底本同处亦印作 $T\mapsto Q(t)$；$T$ 当为参数（或参数空间 $T$ 中的参数）$t$ 之误，正文按语义写作 $t\mapsto Q(t)$。]

% [块边界：原书 p.169（页图 p180）首行为上句的续句“于是可知，Pic Y_n→Pic Z_n(0) 是满射，故只需证 Pic Z_n(0)→Pic Γ 是满射”，归 ch04d；
% 定理 4.39 的证明在 ch04d 中以 [证明（续）] 接续。本块因块边界闭合 proof，会多印一个 ■（原书此处无）。]
